\documentclass[11pt]{article}

\usepackage[margin=1.15in]{geometry}
\usepackage{amsmath,amssymb}
\usepackage{booktabs}
\usepackage{graphicx}
\usepackage{caption}
\usepackage{microtype}
\usepackage[hidelinks]{hyperref}
\usepackage{enumitem}
\usepackage[numbers]{natbib}

\captionsetup{font=small, labelfont=bf, skip=8pt}
\setlist{itemsep=2pt, topsep=4pt}
\newcommand{\code}[1]{\texttt{\small #1}}

\title{\bfseries Constrained Decoding Prevents Cardinality Degradation\\[2pt]
in Structured Extraction}

\author{Raunak Sood\\
\normalsize University of Southern California\\
\normalsize \texttt{raunakso@usc.edu}}

\date{September 2026}

\begin{document}
\maketitle

\begin{abstract}
\noindent
Schema-constrained decoding is widely believed to cost accuracy, because forcing
output through a grammar masks tokens the model would otherwise emit. We measure
that cost along an axis it has not been measured on: the number of items the
model must produce. Holding input length constant by construction, we build
documents containing exactly $k$ verbatim-validated sentences, each naming one
entity, and vary $k$ from 1 to 16. Extraction recall declines in $k$ under
unconstrained decoding, by 0.214 from $k \le 3$ to $k \ge 16$ (Spearman
$-0.293$, $p = 0.011$), and does not decline at all under a JSON schema
($-0.114$, i.e.\ recall rises; $p = 0.73$). Paired on the same 75 documents per
arm, unconstrained decoding leads by 0.200 at $k = 1$ and trails by 0.128 at
$k = 16$ ($p = 0.005$). The mechanism is visible in emission counts: at $k = 16$
the constrained arm emits 9.4 of 16 available items and the unconstrained arm
7.2, at equal precision. Free to choose when to stop, the model stops early, and
stops earlier as the list lengthens; a schema removes that choice because a
partially filled array does not parse. We also show that item count, item
density and document length cannot be held fixed simultaneously, which
determines whether any measurement on this axis is an upper or a lower bound.
\end{abstract}

\section{Introduction}

Ask a language model for a list and it will sometimes give you a shorter one than
the document supports. Whether that tendency worsens as the list grows is an
empirical question with a measurable answer, and it has not been measured in the
setting where it matters most for production systems: structured extraction under
a schema.

Two literatures bracket the question without answering it. Long-context work
measures how retrieval degrades with input length and with the position of the
evidence, and multi-needle variants report that retrieval falls as the number of
inserted facts rises~\citep{multineedle}. Structured-output work measures the cost
of constrained decoding on validity and on accuracy, and reports real
losses~\citep{jsonschemabench}. Neither varies the number of items to emit while
holding input length fixed, and neither compares decoding regimes on that axis.

The distinction matters because the two candidate explanations make opposite
predictions. If degradation is a property of the model's attention over a long
list of targets, a grammar should not help. If it is a property of the model
choosing when to stop, a grammar that rejects incomplete output should help a
great deal. The data support the second.

\paragraph{Contributions.} We fix $k$, the number of items to emit, as an
evaluation axis and hold input length constant by constructing documents rather
than annotating them, which makes $k$ exact and the gold set free. The result is
reported as a slope, not an existence claim: the effect is present under
unconstrained decoding and absent under a schema, inverting the expectation that
grammars cost accuracy. A constraint that any study on this axis inherits is made
explicit. Count, density and length are mechanically linked, so fixing two
determines the third and decides the direction of the residual bias.

\section{Constructing documents with known $k$}

Annotation is the obvious approach and the wrong one. If $k$ is a property of
whichever documents happened to be sampled and a human must recover it, $k$ is
measured with error that correlates with the outcome, since an annotator's
attention flags on exactly the long enumerations where the effect should live.

We construct instead. The item bank is drawn from SEC 10-K filings: sentences
each naming exactly one organisation, each verified to appear verbatim in its
source document. Filler is drawn from the same corpus with every sentence
containing a capitalised run removed, so it cannot carry a legitimate item. A
document is $k$ item sentences injected at uniformly random positions into
filler, padded to a constant 9{,}000 characters. Gold is the injected set.

Three properties follow: $k$ is exact, input length is constant by construction
rather than merely uncorrelated, and position is averaged over rather than
confounded with $k$.

\paragraph{The item bank is screened by measurement, not by rule.} Three
syntactic filters for ``which sentences transplant into a synthetic document''
were tried and all three failed, leaving recall at $k = 1$ far below any usable
baseline. The bank is therefore defined empirically: a candidate sentence is
retained only if the extractor recovers it at $k = 1$ in a short document. 143
of 240 candidates survive.

\paragraph{A validity gate.} Before any slope is fitted, recall at $k = 1$ must
clear 0.80 for a cell to be interpretable: one item, one document, a sentence the
extractor itself produced in the wild. This gate failed three times during
construction, and each failure was informative. The decisive diagnosis came from
holding the item and $k$ fixed and varying only the filler:

\begin{center}
\small
\begin{tabular}{lc}
\toprule
filler & recall at $k=1$ \\
\midrule
1{,}500 characters & 0.667 \\
4{,}000 characters & 0.583 \\
14{,}000 characters & 0.417 \\
\bottomrule
\end{tabular}
\end{center}

Dilution alone is worth 0.25 recall. The filler had been selected \emph{for
containing no entity names}, which made it far more inert than any real prompt,
where a passage selector retains only relationship-dense text. Document length
was reduced accordingly.

\section{The count--density--length trilemma}

The dilution measurement exposes a constraint that binds every study of this
axis. Item density is mechanically determined by the other two quantities:
$\text{density} = (k \times \text{span length}) / \text{document length}$. Fixing
any two fixes the third.

\begin{center}
\small
\begin{tabular}{lll}
\toprule
Held fixed & Varies & Residual confound \\
\midrule
length and count (this work) & density & density aids high $k$; effect is a \textbf{lower bound} \\
density and count & length & length suppresses high $k$; an \textbf{upper bound} \\
length and density & --- & $k$ cannot vary \\
\bottomrule
\end{tabular}
\end{center}

We hold length fixed, so item density rises with $k$: a $k = 16$ document is
denser in items than a $k = 1$ document of the same length. Since dilution
demonstrably suppresses recall, high-$k$ documents receive a density advantage
that works against the hypothesis. Our estimate is therefore conservative.

\section{Design and pre-registration}

The hypothesis, estimator and rejection rule were fixed before any cell was run.
Rejection requires \emph{both} a recall drop of at least 0.15 from $k \le 3$ to
$k \ge 16$ \emph{and} $p < 0.05$ on the rank test. Bucket means are for display;
estimation is per document.

Both arms use \code{qwen3:14b} at temperature 0 with a 32{,}768-token context and
a 24{,}000-token generation ceiling, on 75 matched documents each at
$k \in \{1, 4, 16\}$, 25 per level. The unconstrained arm is prompted for JSON and
parsed post hoc, with a deliberately tolerant parser: a parse failure would score
as a recall failure and confound the arm under test.

\paragraph{The ceiling control matters and is easy to get wrong.} An earlier
attempt used a 30B mixture-of-experts model, which writes 7{,}000 to 12{,}000
tokens unconstrained against an 8{,}192-token window. Four of six $k = 16$
documents truncated mid-array, parsed to zero items, and would have reported
constrained decoding as hugely superior at the level that decides the study. That arm was discarded rather than repaired. In the reported runs, zero
documents truncate and the maximum generation is 686 tokens.

\section{Results}

\begin{figure}[h]
\centering
\includegraphics[width=\textwidth]{figures/cardinality.pdf}
\caption{(a) Recall against item count for both decoding regimes, 75 matched
documents per arm, error bars one standard error. The arms cross between $k=1$
and $k=4$. (b) Items emitted against items available. The constrained arm tracks
availability more closely as $k$ grows; the unconstrained arm falls away.}
\label{fig:main}
\end{figure}

\begin{table}[h]
\centering
\small
\begin{tabular}{lrrrr}
\toprule
$k$ & constrained & unconstrained & paired difference & $p$ \\
\midrule
1 & 0.440 & 0.640 & $+0.200$ & 0.096 \\
4 & 0.640 & 0.530 & $-0.110$ & 0.110 \\
16 & 0.554 & 0.426 & $\mathbf{-0.128}$ & $\mathbf{0.005}$ \\
\bottomrule
\end{tabular}
\caption{Recall by item count and decoding regime, paired on identical documents.}
\label{tab:main}
\end{table}

The pre-registered test rejects for one arm and not the other. Unconstrained
decoding drops $0.214$ from $k \le 3$ to $k \ge 16$ with Spearman $-0.293$,
$p = 0.011$, satisfying both conditions. Constrained decoding drops $-0.114$,
meaning recall rises, with Spearman $+0.041$, $p = 0.73$, satisfying neither.

The $k = 1$ comparison favours unconstrained decoding by 0.200 at $p = 0.096$.
Per-document recall at $k = 1$ is close to a Bernoulli draw, and the variance
arithmetic implies 106 documents are needed to detect a 0.15 difference at 80\%
power against the 25 run. We report it and do not interpret it.

\paragraph{Mechanism.} At $k = 16$ the constrained arm emits 9.4 of 16 available
items and the unconstrained arm 7.2, at precision 0.946 and 0.948 respectively.
The difference is not a precision--recall trade: the unconstrained arm simply
produces fewer items, and produces relatively fewer as the list lengthens. A
schema makes an incomplete array unparseable, so generation continues past the
point where the model would otherwise stop.

\paragraph{Cost.} Unconstrained decoding is not cheaper. It costs $1.1\times$ the
tokens at $k = 1$ and $0.7\times$ at $k = 16$: on this model it is both less
accurate at high $k$ and no less expensive.

\section{Scope conditions}

One model, dense, one family; nothing here separates architecture from decoding
regime. The reasoning-mode flag on this stack changes only which response field
the same tokens arrive in, verified by identical evaluation counts, so this work
says nothing about inference-time reasoning. The item bank is drawn from a single
domain and screened by one extractor, so absolute recall levels are not
comparable across corpora. The slope is the estimand. The $k = 1$ arm is
underpowered by the study's own arithmetic.

\section{Conclusion}

On the axis of how many items a model must emit, the grammar is not what costs
recall. Schema-constrained decoding holds recall flat across $k$, and removing it
produces a decline that clears a pre-registered threshold. The mechanism is a
stopping decision rather than a retrieval failure, which is why removing the
decision fixes it.

For practitioners the implication is direct: free-form generation with post-hoc
parsing, often proposed as an escape from grammar masking, trades away recall on
the long-list cases where extraction is most valuable, and saves nothing.
For evaluation, a single $F_1$ over a benchmark whose $k$ distribution is
unstated is not comparable across benchmarks, because it averages over a mix
nobody reported.

\paragraph{Reproduction.} Code, pre-registration, the constructed evaluation set
and all 450 raw model responses are released at
\url{https://github.com/quantraunak/cardinality-eval}. The reported result
re-scores without a GPU.

\begin{thebibliography}{9}

\bibitem[LangChain(2024)]{multineedle}
LangChain.
\newblock Multi needle in a haystack.
\newblock \url{https://www.langchain.com/blog/multi-needle-in-a-haystack}, 2024.

\bibitem[Geng et~al.(2025)]{jsonschemabench}
S.~Geng et~al.
\newblock JSONSchemaBench: A rigorous benchmark of structured outputs for
  language models.
\newblock arXiv:2501.10868, 2025.

\end{thebibliography}

\end{document}
